
\documentclass[11pt,a4paper]{article}

% ============================================================
% COURS PASSERELLE PHYSIQUE-CHIMIE
% DE LA PREMIÈRE SPÉCIALITÉ À LA TERMINALE SPÉCIALITÉ
% Version rigoureuse, structurée, prête à compiler
% ============================================================

\usepackage[utf8]{inputenc}
\usepackage[T1]{fontenc}
\usepackage[french]{babel}
\usepackage{lmodern}
\usepackage{geometry}
\geometry{margin=1.65cm}

\usepackage{amsmath,amssymb,amsthm}
\usepackage{siunitx}
\usepackage{physics}
\usepackage{array}
\usepackage{tabularx}
\usepackage{longtable}
\usepackage{booktabs}
\usepackage{multicol}
\usepackage{xcolor}
\usepackage{enumitem}
\usepackage{fancyhdr}
\usepackage{hyperref}
\usepackage[most]{tcolorbox}
\usepackage{pifont}
\usepackage{tikz}

\hypersetup{
  colorlinks=true,
  linkcolor=blue!60!black,
  urlcolor=blue!60!black
}

\pagestyle{fancy}
\fancyhf{}
\lhead{Cours passerelle Physique-Chimie}
\rhead{Première $\rightarrow$ Terminale}
\cfoot{\thepage}

\setlist[itemize]{leftmargin=1.25em}
\setlist[enumerate]{leftmargin=1.55em}

\definecolor{bleuprepa}{RGB}{20,55,110}
\definecolor{vertprepa}{RGB}{20,105,70}
\definecolor{rougeprepa}{RGB}{160,35,35}
\definecolor{orangeprepa}{RGB}{180,95,20}
\definecolor{violetprepa}{RGB}{95,45,130}
\definecolor{grisclair}{RGB}{245,247,250}

\newcommand{\casevide}{\ding{113}}
\newcommand{\important}[1]{\textbf{\textcolor{rougeprepa}{#1}}}
\newcommand{\R}{\mathbb{R}}
\newcommand{\N}{\mathbb{N}}
\newcommand{\e}{\mathrm{e}}
\newcommand{\dd}{\,\mathrm{d}}
\newcommand{\vect}[1]{\overrightarrow{#1}}

\sisetup{
  locale = FR,
  detect-all,
  per-mode=symbol
}

\tcbset{
  enhanced,
  breakable,
  sharp corners=downhill,
  boxrule=0.7pt,
  colback=white,
  fonttitle=\bfseries
}

\newtcolorbox{objectif}[1][]{
  colframe=bleuprepa,
  colback=blue!2!white,
  title={Objectif #1}
}

\newtcolorbox{definition}[1][]{
  colframe=bleuprepa,
  colback=cyan!2!white,
  title={Définition #1}
}

\newtcolorbox{propriete}[1][]{
  colframe=vertprepa,
  colback=green!2!white,
  title={Propriété #1}
}

\newtcolorbox{methode}[1][]{
  colframe=violetprepa,
  colback=violet!3!white,
  title={Méthode #1}
}

\newtcolorbox{alerte}[1][]{
  colframe=rougeprepa,
  colback=red!2!white,
  title={Alerte #1}
}

\newtcolorbox{exemple}[1][]{
  colframe=orangeprepa,
  colback=orange!3!white,
  title={Exemple #1}
}

\newtcolorbox{transition}[1][]{
  colframe=black!70,
  colback=grisclair,
  title={Passage première $\rightarrow$ terminale #1}
}

\newtcolorbox{exercice}[1][]{
  colframe=bleuprepa,
  colback=white,
  title={Exercice guidé #1}
}

\title{
  \Huge\textbf{Cours passerelle de Physique-Chimie}\\[2mm]
  \Large De la première spécialité à la terminale spécialité\\[1mm]
  \large Cours rigoureux, méthodes, démonstrations, exercices guidés
}
\author{Document professeur — Lycée général}
\date{}

\begin{document}

\maketitle

\tableofcontents
\newpage

% ============================================================
\section{Introduction : ce qui change entre la première et la terminale}
% ============================================================

\begin{objectif}
En première, on apprend à reconnaître des lois et à les appliquer.  
En terminale, on apprend davantage à \important{modéliser}, \important{interpréter}, \important{résoudre} et \important{critiquer} un modèle.
\end{objectif}

\subsection{Les quatre grands thèmes}

La physique-chimie de spécialité s'organise autour de quatre domaines :

\begin{enumerate}
  \item \textbf{Constitution et transformations de la matière} : transformations chimiques, quantités de matière, dosages, acides-bases, oxydoréduction, cinétique.
  \item \textbf{Mouvement et interactions} : forces, mouvement, lois de Newton, champs, dynamique.
  \item \textbf{L'énergie : conversions et transferts} : énergie mécanique, puissance, bilans énergétiques, transferts thermiques, électricité.
  \item \textbf{Ondes et signaux} : ondes mécaniques, lumière, diffraction, interférences, signaux électriques.
\end{enumerate}

\begin{transition}
Le passage en terminale consiste à transformer les connaissances de première en outils plus généraux :
\[
\text{loi simple} \quad \longrightarrow \quad \text{modèle mathématique}.
\]
Par exemple :
\[
v=\frac{d}{\Delta t}
\quad \longrightarrow \quad
\vec v(t)=\frac{\dd \vec{OM}}{\dd t}.
\]
Ou encore :
\[
n=\frac{m}{M}
\quad \longrightarrow \quad
\text{tableau d'avancement et réactif limitant}.
\]
\end{transition}

\subsection{Les compétences à acquérir}

\begin{multicols}{2}
\begin{itemize}
  \item Identifier le système étudié.
  \item Choisir un modèle.
  \item Écrire une loi physique.
  \item Vérifier l'homogénéité.
  \item Utiliser les unités SI.
  \item Exploiter un graphe.
  \item Linéariser une relation.
  \item Résoudre une équation simple.
  \item Interpréter un résultat.
  \item Critiquer un modèle.
  \item Rédiger une conclusion scientifique.
  \item Passer d'une mesure à une grandeur.
\end{itemize}
\end{multicols}

% ============================================================
\section{Outils fondamentaux : grandeurs, unités et dimensions}
% ============================================================

\subsection{Grandeur physique}

\begin{definition}
Une \textbf{grandeur physique} est une propriété mesurable d'un système. Elle s'exprime par :
\[
\text{valeur numérique} \times \text{unité}.
\]
Par exemple :
\[
m=\SI{2.5}{kg}, \qquad v=\SI{12}{m.s^{-1}}, \qquad E=\SI{150}{J}.
\]
\end{definition}

\begin{alerte}
Une valeur sans unité est généralement inutilisable en physique.  
Écrire seulement
\[
v=12
\]
ne signifie rien si l'on ne précise pas :
\[
v=\SI{12}{m.s^{-1}}
\quad \text{ou} \quad
v=\SI{12}{km.h^{-1}}.
\]
\end{alerte}

\subsection{Unités du Système international}

\begin{center}
\renewcommand{\arraystretch}{1.3}
\begin{tabularx}{\textwidth}{|X|c|c|}
\hline
\textbf{Grandeur} & \textbf{Symbole} & \textbf{Unité SI} \\
\hline
Longueur & $L$ & mètre, \si{m} \\
\hline
Masse & $m$ & kilogramme, \si{kg} \\
\hline
Temps & $t$ & seconde, \si{s} \\
\hline
Température & $T$ & kelvin, \si{K} \\
\hline
Quantité de matière & $n$ & mole, \si{mol} \\
\hline
Intensité électrique & $I$ & ampère, \si{A} \\
\hline
\end{tabularx}
\end{center}

\subsection{Analyse dimensionnelle}

\begin{definition}
La \textbf{dimension} d'une grandeur indique sa nature physique indépendamment de l'unité choisie.  
On note par exemple :
\[
[v]=L\,T^{-1},
\]
ce qui signifie qu'une vitesse est une longueur divisée par un temps.
\end{definition}

\begin{exemple}[énergie cinétique]
L'énergie cinétique est :
\[
E_c=\frac12mv^2.
\]
Dimensions :
\[
[E_c]=[m][v]^2=M(LT^{-1})^2=ML^2T^{-2}.
\]
L'unité correspondante est :
\[
\si{kg.m^2.s^{-2}}=\si{J}.
\]
\end{exemple}

\begin{methode}[vérifier l'homogénéité]
Une formule physique correcte doit être homogène :
\[
\text{dimension du membre gauche}=\text{dimension du membre droit}.
\]
Exemple :
\[
d=vt.
\]
On vérifie :
\[
[d]=L,
\qquad
[vt]=LT^{-1}\times T=L.
\]
La formule est homogène.
\end{methode}

\begin{alerte}
Une formule homogène n'est pas forcément vraie.  
Mais une formule non homogène est forcément fausse.
\end{alerte}

% ============================================================
\section{Chimie I : de la quantité de matière au tableau d'avancement}
% ============================================================

\subsection{Quantité de matière}

\begin{definition}
La quantité de matière $n$ mesure le nombre d'entités chimiques contenues dans un échantillon.  
Elle s'exprime en mole :
\[
n=\frac{N}{N_A},
\]
où $N$ est le nombre d'entités et $N_A$ la constante d'Avogadro :
\[
N_A \approx \SI{6.02e23}{mol^{-1}}.
\]
\end{definition}

\begin{propriete}
Pour un échantillon de masse $m$ et de masse molaire $M$ :
\[
n=\frac{m}{M}.
\]
Avec :
\[
m \text{ en g}, \qquad M \text{ en g.mol}^{-1}, \qquad n \text{ en mol}.
\]
\end{propriete}

\begin{propriete}
Pour une solution de concentration molaire $C$ et de volume $V$ :
\[
n=CV.
\]
Attention :
\[
C \text{ en mol.L}^{-1}, \qquad V \text{ en L}.
\]
\end{propriete}

\begin{alerte}[conversion des volumes]
En chimie au lycée, on utilise souvent :
\[
\SI{1}{mL}=\SI{1e-3}{L}.
\]
Ainsi :
\[
\SI{25.0}{mL}=\SI{2.50e-2}{L}.
\]
\end{alerte}

\subsection{Transformation chimique}

\begin{definition}
Une \textbf{transformation chimique} est une évolution au cours de laquelle des espèces chimiques sont consommées et d'autres formées.
\end{definition}

Une transformation chimique est modélisée par une équation de réaction :
\[
aA+bB \longrightarrow cC+dD.
\]

Les nombres $a,b,c,d$ sont les coefficients stœchiométriques.

\begin{alerte}
Une équation chimique doit respecter :
\begin{itemize}
  \item la conservation des éléments chimiques ;
  \item la conservation de la charge électrique.
\end{itemize}
\end{alerte}

\subsection{Avancement}

\begin{definition}
L'\textbf{avancement} $x$ d'une réaction est une grandeur exprimée en mole qui mesure la progression d'une transformation chimique.
\end{definition}

Pour la réaction :
\[
aA+bB \longrightarrow cC+dD,
\]
si l'avancement est $x$, alors les quantités de matière évoluent ainsi :

\[
n_A=n_{A,0}-ax,
\]
\[
n_B=n_{B,0}-bx,
\]
\[
n_C=n_{C,0}+cx,
\]
\[
n_D=n_{D,0}+dx.
\]

\begin{methode}[tableau d'avancement]
Pour construire un tableau d'avancement :
\begin{enumerate}
  \item écrire l'équation équilibrée ;
  \item écrire les quantités initiales ;
  \item introduire l'avancement $x$ ;
  \item exprimer les quantités finales ;
  \item imposer que les quantités de matière restent positives.
\end{enumerate}
\end{methode}

\begin{exemple}[réaction entre dihydrogène et dioxygène]
On considère :
\[
2H_2+O_2\longrightarrow 2H_2O.
\]

\[
\begin{array}{c|ccc}
 & H_2 & O_2 & H_2O \\
\hline
\text{Initial} & n_{H_2,0} & n_{O_2,0} & 0 \\
\text{Évolution} & -2x & -x & +2x \\
\text{Final} & n_{H_2,0}-2x & n_{O_2,0}-x & 2x \\
\end{array}
\]
\end{exemple}

\subsection{Réactif limitant}

\begin{definition}
Le \textbf{réactif limitant} est le réactif totalement consommé en premier. Il impose l'avancement maximal $x_{\max}$.
\end{definition}

Pour la réaction :
\[
aA+bB\longrightarrow cC+dD,
\]
on a :
\[
x\leq \frac{n_{A,0}}{a},
\qquad
x\leq \frac{n_{B,0}}{b}.
\]
Donc :
\[
x_{\max}=\min\left(\frac{n_{A,0}}{a},\frac{n_{B,0}}{b}\right).
\]

\begin{transition}
En première, on apprend à calculer un réactif limitant.  
En terminale, cette idée devient centrale pour :
\begin{itemize}
  \item les dosages ;
  \item les transformations non totales ;
  \item les équilibres chimiques ;
  \item les rendements de synthèse ;
  \item la cinétique.
\end{itemize}
\end{transition}

% ============================================================
\section{Chimie II : des dosages aux titrages}
% ============================================================

\subsection{Dosage}

\begin{definition}
Un \textbf{dosage} consiste à déterminer la concentration d'une espèce chimique dans une solution.
\end{definition}

Il existe plusieurs types de dosages :
\begin{itemize}
  \item dosage par étalonnage ;
  \item titrage direct ;
  \item titrage pH-métrique ;
  \item titrage conductimétrique ;
  \item titrage colorimétrique.
\end{itemize}

\subsection{Titrage}

\begin{definition}
Un \textbf{titrage} est un dosage reposant sur une réaction chimique entre :
\begin{itemize}
  \item l'espèce titrée, de concentration inconnue ;
  \item l'espèce titrante, de concentration connue.
\end{itemize}
\end{definition}

\subsection{Équivalence}

\begin{definition}
L'\textbf{équivalence} d'un titrage est l'instant où les réactifs titré et titrant ont été introduits dans les proportions stœchiométriques.
\end{definition}

Pour une réaction :
\[
aA+bB\longrightarrow \text{produits},
\]
si $A$ est titré et $B$ est titrant, alors à l'équivalence :
\[
\frac{n_A}{a}=\frac{n_B}{b}.
\]

\begin{methode}[exploiter un titrage]
\begin{enumerate}
  \item Identifier l'espèce titrée et l'espèce titrante.
  \item Écrire l'équation support du titrage.
  \item Repérer l'équivalence.
  \item Utiliser la relation stœchiométrique.
  \item Remplacer par $n=CV$.
  \item Isoler la concentration inconnue.
  \item Vérifier les unités.
\end{enumerate}
\end{methode}

\begin{exemple}[titrage acide-base simple]
On dose un acide $AH$ par une base $HO^-$.  
Réaction :
\[
AH+HO^- \longrightarrow A^-+H_2O.
\]
Les coefficients sont tous égaux à 1.  
À l'équivalence :
\[
n(AH)=n(HO^-).
\]
Donc :
\[
C_A V_A=C_B V_E.
\]
Ainsi :
\[
C_A=\frac{C_BV_E}{V_A}.
\]
\end{exemple}

\begin{alerte}
Les volumes doivent être dans la même unité.  
Si $V_A$ est en L et $V_E$ en mL, la formule devient numériquement fausse.
\end{alerte}

\subsection{Passage vers la terminale}

\begin{transition}
En première, le titrage est souvent vu comme une application de $n=CV$.  
En terminale, il faut aussi comprendre :
\begin{itemize}
  \item pourquoi l'équivalence se repère par un saut de pH ou une rupture de pente ;
  \item comment choisir un indicateur coloré ;
  \item comment exploiter une courbe expérimentale ;
  \item comment discuter les incertitudes.
\end{itemize}
\end{transition}

% ============================================================
\section{Chimie III : acides, bases et pH}
% ============================================================

\subsection{Définition de Brønsted}

\begin{definition}
Un \textbf{acide} au sens de Brønsted est une espèce capable de céder un proton $H^+$.  
Une \textbf{base} au sens de Brønsted est une espèce capable de capter un proton $H^+$.
\end{definition}

Exemple :
\[
AH \rightleftharpoons A^-+H^+.
\]
Ici :
\[
AH/A^-
\]
est un couple acide-base.

\subsection{pH}

\begin{definition}
Le pH d'une solution aqueuse est défini par :
\[
pH=-\log[H_3O^+].
\]
La concentration $[H_3O^+]$ est exprimée en \si{mol.L^{-1}}.
\end{definition}

On en déduit :
\[
[H_3O^+]=10^{-pH}.
\]

\begin{exemple}
Si :
\[
pH=3,
\]
alors :
\[
[H_3O^+]=10^{-3}\ \si{mol.L^{-1}}.
\]
\end{exemple}

\begin{alerte}
Le pH est une grandeur logarithmique.  
Une baisse d'une unité de pH signifie que la concentration en ions oxonium est multipliée par 10.
\end{alerte}

\subsection{Acide fort et acide faible}

\begin{definition}
Un acide fort réagit totalement avec l'eau.  
Un acide faible réagit partiellement avec l'eau.
\end{definition}

Acide fort :
\[
AH+H_2O \longrightarrow A^-+H_3O^+.
\]

Acide faible :
\[
AH+H_2O \rightleftharpoons A^-+H_3O^+.
\]

\begin{transition}
En première, on manipule surtout des réactions considérées comme totales.  
En terminale, on rencontre des réactions non totales et des équilibres chimiques. Le symbole :
\[
\rightleftharpoons
\]
devient fondamental.
\end{transition}

% ============================================================
\section{Chimie IV : oxydoréduction}
% ============================================================

\subsection{Oxydant et réducteur}

\begin{definition}
Un \textbf{oxydant} est une espèce capable de capter des électrons.  
Un \textbf{réducteur} est une espèce capable de céder des électrons.
\end{definition}

Un couple oxydant/réducteur s'écrit :
\[
Ox/Red.
\]

La demi-équation associée est :
\[
Ox+ne^- \rightleftharpoons Red.
\]

\subsection{Équilibrer une demi-équation}

\begin{methode}[milieu acide]
Pour équilibrer une demi-équation en milieu acide :
\begin{enumerate}
  \item équilibrer les atomes autres que $O$ et $H$ ;
  \item équilibrer les atomes d'oxygène avec $H_2O$ ;
  \item équilibrer les atomes d'hydrogène avec $H^+$ ;
  \item équilibrer les charges avec $e^-$ ;
  \item vérifier la conservation des charges et des atomes.
\end{enumerate}
\end{methode}

\begin{exemple}
Pour le couple :
\[
MnO_4^-/Mn^{2+},
\]
en milieu acide :
\[
MnO_4^-+8H^++5e^- \longrightarrow Mn^{2+}+4H_2O.
\]
\end{exemple}

\subsection{Réaction d'oxydoréduction}

Une réaction d'oxydoréduction est un transfert d'électrons entre un oxydant et un réducteur.

\begin{exemple}
Avec les couples :
\[
Cu^{2+}/Cu
\qquad \text{et} \qquad
Zn^{2+}/Zn,
\]
on peut écrire :
\[
Cu^{2+}+2e^- \longrightarrow Cu,
\]
\[
Zn \longrightarrow Zn^{2+}+2e^-.
\]
En additionnant :
\[
Cu^{2+}+Zn \longrightarrow Cu+Zn^{2+}.
\]
\end{exemple}

\begin{transition}
En terminale, l'oxydoréduction intervient dans :
\begin{itemize}
  \item les piles ;
  \item les électrolyses ;
  \item les titrages redox ;
  \item les transferts d'énergie chimique vers énergie électrique.
\end{itemize}
\end{transition}

% ============================================================
\section{Mécanique I : mouvement, vitesse et accélération}
% ============================================================

\subsection{Système et référentiel}

\begin{definition}
Le \textbf{système} est l'objet ou l'ensemble d'objets que l'on étudie.
\end{definition}

\begin{definition}
Un \textbf{référentiel} est un solide de référence par rapport auquel on décrit le mouvement.
\end{definition}

\begin{alerte}
Avant toute application d'une loi de mécanique, il faut écrire :
\[
\text{Système étudié : ...}
\]
\[
\text{Référentiel choisi : ...}
\]
\]
Sans cela, la rédaction est incomplète.
\end{alerte}

\subsection{Position}

Dans un repère, la position d'un point matériel $M$ est donnée par son vecteur position :
\[
\vec{OM}(t).
\]

En coordonnées cartésiennes :
\[
\vec{OM}(t)=x(t)\vec i+y(t)\vec j+z(t)\vec k.
\]

\subsection{Vitesse}

\begin{definition}
Le vecteur vitesse est la dérivée temporelle du vecteur position :
\[
\vec v(t)=\frac{\dd \vec{OM}}{\dd t}.
\]
\end{definition}

En coordonnées :
\[
\vec v(t)=x'(t)\vec i+y'(t)\vec j+z'(t)\vec k.
\]

\subsection{Accélération}

\begin{definition}
Le vecteur accélération est la dérivée temporelle du vecteur vitesse :
\[
\vec a(t)=\frac{\dd \vec v}{\dd t}.
\]
\end{definition}

Donc :
\[
\vec a(t)=\frac{\dd^2\vec{OM}}{\dd t^2}.
\]

En coordonnées :
\[
\vec a(t)=x''(t)\vec i+y''(t)\vec j+z''(t)\vec k.
\]

\begin{transition}
En première, on calcule souvent une vitesse moyenne :
\[
v=\frac{d}{\Delta t}.
\]
En terminale, on utilise des grandeurs instantanées :
\[
\vec v(t)=\frac{\dd \vec{OM}}{\dd t},
\qquad
\vec a(t)=\frac{\dd \vec v}{\dd t}.
\]
\end{transition}

% ============================================================
\section{Mécanique II : forces et lois de Newton}
% ============================================================

\subsection{Force}

\begin{definition}
Une force modélise une action mécanique exercée par un objet sur un autre.  
Elle est représentée par un vecteur :
\[
\vec F.
\]
Son unité est le newton :
\[
\si{N}.
\]
\end{definition}

\subsection{Poids}

\begin{definition}
Le poids d'un objet de masse $m$ au voisinage de la Terre est :
\[
\vec P=m\vec g.
\]
\end{definition}

Avec :
\[
g\simeq \SI{9.81}{m.s^{-2}}.
\]

\subsection{Deuxième loi de Newton}

\begin{propriete}
Dans un référentiel galiléen, pour un système de masse constante :
\[
\sum \vec F_{\text{ext}}=m\vec a.
\]
\end{propriete}

\begin{methode}[appliquer la deuxième loi de Newton]
\begin{enumerate}
  \item Définir le système.
  \item Choisir un référentiel supposé galiléen.
  \item Faire le bilan des forces extérieures.
  \item Écrire :
  \[
  \sum \vec F_{\text{ext}}=m\vec a.
  \]
  \item Choisir un repère.
  \item Projeter l'équation vectorielle.
  \item Résoudre.
  \item Interpréter.
\end{enumerate}
\end{methode}

\subsection{Chute libre verticale}

On considère un objet lancé verticalement sans frottement.  
Le système est l'objet.  
La seule force est le poids :
\[
\vec P=m\vec g.
\]

Deuxième loi de Newton :
\[
m\vec a=m\vec g.
\]
Donc :
\[
\vec a=\vec g.
\]

Si l'axe vertical est orienté vers le haut :
\[
a_y=-g.
\]

On intègre :
\[
v_y(t)=v_0-gt.
\]

Puis :
\[
y(t)=y_0+v_0t-\frac12gt^2.
\]

\begin{transition}
En première, on utilise parfois directement des formules de mouvement.  
En terminale, il faut savoir les retrouver à partir de Newton.
\end{transition}

% ============================================================
\section{Énergie mécanique}
% ============================================================

\subsection{Énergie cinétique}

\begin{definition}
L'énergie cinétique d'un point matériel de masse $m$ et de vitesse $v$ est :
\[
E_c=\frac12mv^2.
\]
\end{definition}

\subsection{Énergie potentielle de pesanteur}

\begin{definition}
Dans un champ de pesanteur uniforme, l'énergie potentielle de pesanteur peut s'écrire :
\[
E_p=mgz+C.
\]
En choisissant une origine des énergies, on prend souvent :
\[
E_p=mgz.
\]
\end{definition}

\subsection{Énergie mécanique}

\begin{definition}
L'énergie mécanique est :
\[
E_m=E_c+E_p.
\]
\end{definition}

\subsection{Conservation de l'énergie mécanique}

\begin{propriete}
Si seules des forces conservatives travaillent, alors l'énergie mécanique se conserve :
\[
E_m=\text{constante}.
\]
\end{propriete}

\begin{exemple}[chute libre sans frottement]
Si un objet tombe sans frottement :
\[
E_c+E_p=\text{constante}.
\]
Donc :
\[
\frac12mv^2+mgz=\text{constante}.
\]
\end{exemple}

\subsection{Travail d'une force}

\begin{definition}
Pour une force constante $\vec F$ appliquée sur un déplacement $\vec{AB}$ :
\[
W_{A\to B}(\vec F)=\vec F\cdot \vec{AB}.
\]
\end{definition}

Donc :
\[
W_{A\to B}(\vec F)=F\cdot AB\cdot \cos(\theta).
\]

\begin{alerte}
Le travail est positif si la force favorise le mouvement.  
Il est négatif si elle s'oppose au mouvement.
\end{alerte}

\subsection{Théorème de l'énergie cinétique}

\begin{propriete}
La variation d'énergie cinétique d'un système est égale à la somme des travaux des forces extérieures :
\[
\Delta E_c=\sum W(\vec F_{\text{ext}}).
\]
\end{propriete}

\begin{transition}
En première, on calcule l'énergie cinétique et potentielle.  
En terminale, on choisit entre :
\[
\sum \vec F=m\vec a
\]
et
\[
\Delta E_c=\sum W.
\]
Le choix dépend du problème :
\begin{itemize}
  \item Newton donne le mouvement ;
  \item l'énergie donne souvent plus vite les vitesses ou les hauteurs.
\end{itemize}
\end{transition}

% ============================================================
\section{Électricité : de la loi d'Ohm au circuit RC}
% ============================================================

\subsection{Intensité, tension, résistance}

\begin{definition}
L'intensité électrique $I$ est un débit de charge :
\[
I=\frac{\Delta q}{\Delta t}.
\]
Son unité est l'ampère :
\[
\si{A}.
\]
\end{definition}

\begin{definition}
La tension électrique $U$ entre deux points est une différence de potentiel électrique.  
Son unité est le volt :
\[
\si{V}.
\]
\end{definition}

\begin{propriete}[loi d'Ohm]
Pour un conducteur ohmique :
\[
U=RI.
\]
\end{propriete}

\subsection{Puissance électrique}

\begin{propriete}
La puissance électrique reçue par un dipôle est :
\[
P=UI.
\]
Pour un conducteur ohmique :
\[
P=RI^2=\frac{U^2}{R}.
\]
\end{propriete}

\subsection{Énergie électrique}

Si une puissance $P$ est reçue pendant une durée $\Delta t$, l'énergie reçue est :
\[
E=P\Delta t.
\]

\begin{alerte}
En physique :
\[
\SI{1}{W}=\SI{1}{J.s^{-1}}.
\]
Donc :
\[
E=P\Delta t
\]
est en joule si $P$ est en watt et $\Delta t$ en seconde.
\end{alerte}

\subsection{Condensateur}

\begin{definition}
Un condensateur est un dipôle capable de stocker de l'énergie électrique.  
La charge $q$ d'un condensateur est liée à sa tension $u_C$ par :
\[
q=Cu_C.
\]
\end{definition}

$C$ est la capacité du condensateur, en farad :
\[
\si{F}.
\]

\subsection{Circuit RC}

Dans un circuit série composé d'un générateur de tension $E$, d'une résistance $R$ et d'un condensateur $C$, la loi des mailles donne :
\[
E=u_R+u_C.
\]

Or :
\[
u_R=Ri,
\]
et :
\[
i=\frac{\dd q}{\dd t}=C\frac{\dd u_C}{\dd t}.
\]

Donc :
\[
E=RC\frac{\dd u_C}{\dd t}+u_C.
\]

On obtient l'équation différentielle :
\[
\frac{\dd u_C}{\dd t}+\frac{1}{RC}u_C=\frac{E}{RC}.
\]

\begin{definition}
La constante de temps du circuit est :
\[
\tau=RC.
\]
\end{definition}

La solution de charge, si $u_C(0)=0$, est :
\[
u_C(t)=E\left(1-\e^{-t/\tau}\right).
\]

\begin{transition}
En première, on utilise la loi d'Ohm et la puissance électrique.  
En terminale, les circuits deviennent dynamiques : les grandeurs dépendent du temps et vérifient des équations différentielles.
\end{transition}

% ============================================================
\section{Ondes : de la propagation à la diffraction}
% ============================================================

\subsection{Onde mécanique}

\begin{definition}
Une onde mécanique est la propagation d'une perturbation dans un milieu matériel, sans transport global de matière.
\end{definition}

Exemples :
\begin{itemize}
  \item onde sonore ;
  \item vague à la surface de l'eau ;
  \item onde sur une corde.
\end{itemize}

\subsection{Grandeurs associées à une onde périodique}

\begin{definition}
La période $T$ est la durée d'un motif temporel.  
La fréquence $f$ est :
\[
f=\frac{1}{T}.
\]
\end{definition}

\begin{definition}
La longueur d'onde $\lambda$ est la distance parcourue par l'onde pendant une période.
\end{definition}

Si l'onde se propage à la célérité $v$, alors :
\[
v=\lambda f.
\]
Ou encore :
\[
v=\frac{\lambda}{T}.
\]

\subsection{Diffraction}

\begin{definition}
La diffraction est l'étalement d'une onde lorsqu'elle rencontre une ouverture ou un obstacle de dimension comparable à sa longueur d'onde.
\end{definition}

Pour une fente de largeur $a$, l'angle caractéristique de diffraction vérifie approximativement :
\[
\theta\simeq \frac{\lambda}{a}.
\]

\begin{alerte}
La diffraction est d'autant plus marquée que :
\[
a
\]
est petit devant ou comparable à :
\[
\lambda.
\]
\end{alerte}

\subsection{Interférences}

\begin{definition}
Les interférences résultent de la superposition de deux ondes cohérentes.
\end{definition}

Pour deux sources cohérentes séparées par une distance $a$, sur un écran à distance $D$, l'interfrange vaut :
\[
i=\frac{\lambda D}{a}.
\]

\begin{transition}
En première, on étudie les ondes mécaniques et les signaux.  
En terminale, on approfondit les phénomènes ondulatoires :
\begin{itemize}
  \item diffraction ;
  \item interférences ;
  \item mesure de longueurs d'onde ;
  \item exploitation quantitative de figures optiques.
\end{itemize}
\end{transition}

% ============================================================
\section{Mesures, incertitudes et exploitation de données}
% ============================================================

\subsection{Mesure}

\begin{definition}
Une mesure est une estimation expérimentale d'une grandeur.  
Elle n'est jamais parfaitement exacte.
\end{definition}

Une mesure doit être donnée avec :
\[
\text{valeur mesurée} \quad + \quad \text{unité}.
\]

\subsection{Incertitude}

\begin{definition}
L'incertitude mesure la dispersion possible autour d'une valeur mesurée.
\end{definition}

On écrit souvent :
\[
x=x_{\text{mesuré}}\pm u(x).
\]

\begin{exemple}
\[
L=(12{,}4\pm0{,}1)\ \si{cm}.
\]
Cela signifie que la valeur réelle est probablement proche de :
\[
\SI{12.4}{cm}.
\]
\end{exemple}

\subsection{Écart relatif}

\begin{definition}
L'écart relatif entre une valeur expérimentale $x_{\text{exp}}$ et une valeur de référence $x_{\text{ref}}$ est :
\[
\varepsilon=\frac{|x_{\text{exp}}-x_{\text{ref}}|}{x_{\text{ref}}}.
\]
\end{definition}

En pourcentage :
\[
\varepsilon_{\%}=100\times \frac{|x_{\text{exp}}-x_{\text{ref}}|}{x_{\text{ref}}}.
\]

\subsection{Exploitation graphique}

\begin{methode}[exploiter une droite]
Si une relation est de la forme :
\[
y=ax+b,
\]
alors :
\begin{itemize}
  \item $a$ est le coefficient directeur ;
  \item $b$ est l'ordonnée à l'origine.
\end{itemize}
\end{methode}

\begin{exemple}[loi de Beer-Lambert]
La loi de Beer-Lambert s'écrit :
\[
A=kC.
\]
Si l'on trace $A$ en fonction de $C$, on obtient une droite passant par l'origine.  
Le coefficient directeur vaut :
\[
k.
\]
\end{exemple}

% ============================================================
\section{Méthodes transversales de rédaction}
% ============================================================

\begin{methode}[rédiger une question de mécanique]
\begin{enumerate}
  \item Système : préciser l'objet étudié.
  \item Référentiel : préciser le référentiel.
  \item Forces : faire le bilan.
  \item Loi : écrire la loi de Newton ou l'énergie.
  \item Projection : si nécessaire.
  \item Calcul : avec unités.
  \item Conclusion : phrase physique.
\end{enumerate}
\end{methode}

\begin{methode}[rédiger une question de chimie]
\begin{enumerate}
  \item Identifier les espèces chimiques.
  \item Écrire l'équation équilibrée.
  \item Calculer les quantités de matière.
  \item Utiliser la stœchiométrie.
  \item Calculer la grandeur demandée.
  \item Vérifier les unités.
  \item Interpréter.
\end{enumerate}
\end{methode}

\begin{methode}[rédiger une question d'électricité]
\begin{enumerate}
  \item Identifier les dipôles.
  \item Choisir les conventions.
  \item Écrire la loi des mailles ou des nœuds.
  \item Utiliser les lois des dipôles.
  \item Résoudre l'équation.
  \item Vérifier l'homogénéité.
  \item Interpréter la limite ou le régime permanent.
\end{enumerate}
\end{methode}

% ============================================================
\section{Exercices guidés de transition}
% ============================================================

\begin{exercice}[1 — Tableau d'avancement et réactif limitant]
On introduit $n_0(H_2)=\SI{0.30}{mol}$ et $n_0(O_2)=\SI{0.10}{mol}$.  
La réaction est :
\[
2H_2+O_2\longrightarrow 2H_2O.
\]

\begin{enumerate}
  \item Construire le tableau d'avancement.
  \item Déterminer $x_{\max}$.
  \item Identifier le réactif limitant.
  \item Calculer la quantité d'eau formée.
\end{enumerate}

\begin{methode}[Correction]
On a :
\[
x\leq \frac{0{,}30}{2}=0{,}15,
\qquad
x\leq \frac{0{,}10}{1}=0{,}10.
\]
Donc :
\[
x_{\max}=0{,}10\ \si{mol}.
\]
Le dioxygène est limitant.  
La quantité d'eau formée vaut :
\[
n(H_2O)=2x_{\max}=0{,}20\ \si{mol}.
\]
\end{methode}
\end{exercice}

\begin{exercice}[2 — Titrage acide-base]
On dose $V_A=\SI{20.0}{mL}$ d'une solution d'acide chlorhydrique par une solution de soude de concentration :
\[
C_B=\SI{0.100}{mol.L^{-1}}.
\]
L'équivalence est obtenue pour :
\[
V_E=\SI{14.6}{mL}.
\]

\begin{enumerate}
  \item Écrire l'équation de réaction.
  \item Écrire la relation à l'équivalence.
  \item Déterminer la concentration $C_A$.
\end{enumerate}

\begin{methode}[Correction]
Réaction :
\[
H_3O^+ + HO^- \longrightarrow 2H_2O.
\]
À l'équivalence :
\[
n(H_3O^+)=n(HO^-).
\]
Donc :
\[
C_AV_A=C_BV_E.
\]
Ainsi :
\[
C_A=\frac{C_BV_E}{V_A}.
\]
Avec les volumes dans la même unité :
\[
C_A=\frac{0{,}100\times14{,}6}{20{,}0}
=0{,}0730\ \si{mol.L^{-1}}.
\]
\end{methode}
\end{exercice}

\begin{exercice}[3 — Chute verticale sans frottement]
Une balle est lâchée sans vitesse initiale depuis une hauteur :
\[
h=\SI{20}{m}.
\]
On néglige les frottements.

\begin{enumerate}
  \item Déterminer l'accélération.
  \item Écrire l'équation horaire de la hauteur.
  \item Déterminer la durée de chute.
  \item Déterminer la vitesse juste avant l'impact.
\end{enumerate}

\begin{methode}[Correction]
Système : balle.  
Référentiel : terrestre supposé galiléen.  
Force : poids.

Deuxième loi de Newton :
\[
m\vec a=m\vec g.
\]
Donc :
\[
\vec a=\vec g.
\]

Si l'axe est orienté vers le haut :
\[
z(t)=h-\frac12gt^2.
\]
L'impact a lieu quand $z(t)=0$ :
\[
h-\frac12gt^2=0.
\]
Donc :
\[
t=\sqrt{\frac{2h}{g}}.
\]
Numériquement :
\[
t\simeq \sqrt{\frac{40}{9{,}81}}\simeq \SI{2.02}{s}.
\]
La vitesse vaut :
\[
v(t)=-gt.
\]
Donc :
\[
|v|\simeq 9{,}81\times2{,}02\simeq \SI{19.8}{m.s^{-1}}.
\]
\end{methode}
\end{exercice}

\begin{exercice}[4 — Circuit RC]
Un condensateur de capacité :
\[
C=\SI{100}{\micro F}
\]
est chargé à travers une résistance :
\[
R=\SI{10}{k\ohm}.
\]

\begin{enumerate}
  \item Calculer la constante de temps $\tau$.
  \item Écrire l'expression de $u_C(t)$ lors de la charge sous une tension $E$.
  \item Calculer la tension au bout d'une constante de temps.
\end{enumerate}

\begin{methode}[Correction]
\[
\tau=RC.
\]
Avec :
\[
R=\SI{1.0e4}{\ohm},
\qquad
C=\SI{1.0e-4}{F}.
\]
Donc :
\[
\tau=1{,}0\ \si{s}.
\]
Lors de la charge :
\[
u_C(t)=E(1-\e^{-t/\tau}).
\]
Pour $t=\tau$ :
\[
u_C(\tau)=E(1-\e^{-1})\approx0{,}63E.
\]
\end{methode}
\end{exercice}

\begin{exercice}[5 — Diffraction]
Un laser de longueur d'onde :
\[
\lambda=\SI{650}{nm}
\]
éclaire une fente de largeur $a$.  
On observe une tache centrale de largeur :
\[
L=\SI{4.0}{cm}
\]
sur un écran situé à :
\[
D=\SI{2.0}{m}.
\]

\begin{enumerate}
  \item Relier $\theta$, $L$ et $D$.
  \item Utiliser $\theta\simeq \lambda/a$.
  \item Déterminer $a$.
\end{enumerate}

\begin{methode}[Correction]
Pour petits angles :
\[
\theta\simeq \frac{L/2}{D}=\frac{L}{2D}.
\]
Or :
\[
\theta\simeq \frac{\lambda}{a}.
\]
Donc :
\[
a\simeq \frac{2D\lambda}{L}.
\]
Conversions :
\[
\lambda=\SI{650e-9}{m},
\qquad
L=\SI{4.0e-2}{m}.
\]
Donc :
\[
a=\frac{2\times2{,}0\times650\times10^{-9}}{4{,}0\times10^{-2}}
=6{,}5\times10^{-5}\ \si{m}.
\]
Ainsi :
\[
a=\SI{65}{\micro m}.
\]
\end{methode}
\end{exercice}

% ============================================================
\section{Checklist de passage en terminale}
% ============================================================

\begin{objectif}
Avant d'entrer sereinement en terminale spécialité physique-chimie, il faut savoir faire les points suivants sans hésitation.
\end{objectif}

\begin{multicols}{2}
\begin{itemize}
  \item[\casevide] Convertir les unités.
  \item[\casevide] Vérifier l'homogénéité d'une formule.
  \item[\casevide] Calculer une quantité de matière.
  \item[\casevide] Construire un tableau d'avancement.
  \item[\casevide] Identifier un réactif limitant.
  \item[\casevide] Exploiter un titrage.
  \item[\casevide] Définir un acide et une base.
  \item[\casevide] Calculer un pH simple.
  \item[\casevide] Équilibrer une demi-équation redox.
  \item[\casevide] Définir système et référentiel.
  \item[\casevide] Faire un bilan des forces.
  \item[\casevide] Appliquer la deuxième loi de Newton.
  \item[\casevide] Utiliser l'énergie cinétique.
  \item[\casevide] Utiliser l'énergie potentielle.
  \item[\casevide] Exploiter un circuit électrique simple.
  \item[\casevide] Comprendre un circuit RC.
  \item[\casevide] Utiliser $v=\lambda f$.
  \item[\casevide] Utiliser une formule de diffraction.
  \item[\casevide] Exploiter un graphique.
  \item[\casevide] Rédiger une conclusion physique.
\end{itemize}
\end{multicols}

% ============================================================
\section{Conclusion}
% ============================================================

\begin{alerte}[Idée essentielle]
La terminale ne demande pas seulement plus de formules.  
Elle demande surtout une meilleure capacité à passer :
\[
\text{situation réelle}
\quad \longrightarrow \quad
\text{modèle}
\quad \longrightarrow \quad
\text{équation}
\quad \longrightarrow \quad
\text{résultat}
\quad \longrightarrow \quad
\text{interprétation}.
\]
\end{alerte}

\begin{center}
\Large
\textbf{La rigueur en physique-chimie repose sur trois piliers :}\\[1mm]
\textbf{les lois, les unités, et l'interprétation.}
\end{center}

\end{document}
```

[1]: https://eduscol.education.gouv.fr/5829/programmes-et-ressources-en-physique-chimie-voie-gt?utm_source=chatgpt.com "Programmes et ressources en physique-chimie - voie GT"
